% Appendix: Paraunitarity Proofs
% Converted on 2026-10-08 from docs/lswt/04-appendices/paraunitarity-proofs.md (status: draft, source-section: "Paraunitary Diagonalization (source pp. 10–11); restructured TeX appendix 'Para-unitarity proofs'").
% Review status: draft; awaiting the user's physics and mathematics review.

\hypertarget{sec-lswt-paraunitarity-proofs}{%
\section{Appendix: Paraunitarity Proofs}\label{sec-lswt-paraunitarity-proofs}}

The diagonalization of a bosonic Bogoliubov--de Gennes (BdG) matrix rests on a few algebraic facts about the indefinite metric \(\Sigma_3=\sigma_3\otimes\mathsf I_{N_{\mathrm{sub}}}\): paraunitary matrices form a group, the eigenvectors of \(\Sigma_3\mathsf H\) are orthogonal with respect to \(\Sigma_3\), and positive definiteness of \(\mathsf H\) fixes both the reality of the spectrum and the sign of each eigenvector's \(\Sigma_3\)-norm. This appendix proves these facts and the properties of the transformation \(\mathsf T_{\mathbf k}\) used in Paraunitary Diagonalization (Section~\ref{sec-lswt-paraunitary-diagonalization}). Throughout, \(\mathsf H\) is a Hermitian \(2N_{\mathrm{sub}}\times2N_{\mathrm{sub}}\) matrix, and the momentum index is suppressed where only one momentum enters.

\subsection{Paraunitary Matrices Form a Group}

A matrix \(\mathsf T\) is paraunitary when \(\mathsf T\Sigma_3\mathsf T^\dagger=\Sigma_3\). Taking determinants gives \(|\det\mathsf T|^2=1\), so \(\mathsf T\) is invertible. Multiplying the condition by \(\Sigma_3\) on the right gives \(\mathsf T(\Sigma_3\mathsf T^\dagger\Sigma_3)=\mathsf I\), so

\[
\mathsf T^{-1}=\Sigma_3\mathsf T^\dagger\Sigma_3.
\]

Inserting this inverse into \(\mathsf T^{-1}\mathsf T=\mathsf I\) gives \(\Sigma_3\mathsf T^\dagger\Sigma_3\mathsf T=\mathsf I\), which is the second form \(\mathsf T^\dagger\Sigma_3\mathsf T=\Sigma_3\). The argument runs equally in the other direction, so the two forms of the paraunitary condition are equivalent. If \(\mathsf T_1\) and \(\mathsf T_2\) are paraunitary, then \(\mathsf T_1\mathsf T_2\Sigma_3\mathsf T_2^\dagger\mathsf T_1^\dagger=\mathsf T_1\Sigma_3\mathsf T_1^\dagger=\Sigma_3\), and the inverse \(\Sigma_3\mathsf T_1^\dagger\Sigma_3\) satisfies the second form of the condition. Paraunitary matrices therefore form a group, the pseudo-unitary group \(\mathrm U(N_{\mathrm{sub}},N_{\mathrm{sub}})\). A paraunitary matrix that is also unitary commutes with \(\Sigma_3\) and is block diagonal in the particle and hole components.

\subsection{Metric Orthogonality of Eigenvectors}

Let \(\Sigma_3\mathsf H\mathbf v=\lambda\mathbf v\) and \(\Sigma_3\mathsf H\mathbf w=\mu\mathbf w\). Multiplying the first relation by \(\mathbf w^\dagger\Sigma_3\) and the conjugate of the second by \(\mathbf v\), and using \(\Sigma_3^2=\mathsf I\) and \(\mathsf H^\dagger=\mathsf H\), gives

\begin{equation}\label{eq-lswt-metric-orthogonality}
\mathbf w^\dagger\mathsf H\mathbf v=\lambda\,\mathbf w^\dagger\Sigma_3\mathbf v=\mu^*\,\mathbf w^\dagger\Sigma_3\mathbf v,
\qquad
(\lambda-\mu^*)\,\mathbf w^\dagger\Sigma_3\mathbf v=0.
\end{equation}

Eigenvectors whose eigenvalues satisfy \(\lambda\ne\mu^*\) are therefore orthogonal in the indefinite inner product \(\mathbf w^\dagger\Sigma_3\mathbf v\). For \(\mathbf w=\mathbf v\) the relation reads \(\mathbf v^\dagger\mathsf H\mathbf v=\lambda\,\mathbf v^\dagger\Sigma_3\mathbf v\), and an eigenvector with a nonreal eigenvalue has \(\mathbf v^\dagger\Sigma_3\mathbf v=0\).

\subsection{Consequences of Positive Definiteness}

Assume that \(\mathsf H\) is positive definite. Three properties of \(\Sigma_3\mathsf H\) follow.

\textbf{Real, nonzero eigenvalues with a fixed metric sign.} For an eigenvector \(\mathbf v\), the relation \(\mathbf v^\dagger\mathsf H\mathbf v=\lambda\,\mathbf v^\dagger\Sigma_3\mathbf v\) has a positive left-hand side. If \(\lambda\) were not real, the previous section would give \(\mathbf v^\dagger\Sigma_3\mathbf v=0\) and hence \(\mathbf v^\dagger\mathsf H\mathbf v=0\), a contradiction. Therefore \(\lambda\) is real and nonzero, \(\mathbf v^\dagger\Sigma_3\mathbf v\ne0\), and

\begin{equation}\label{eq-lswt-metric-sign-rule}
\operatorname{sign}\lambda=\operatorname{sign}\big(\mathbf v^\dagger\Sigma_3\mathbf v\big).
\end{equation}

Eigenvectors with positive eigenvalues have a positive \(\Sigma_3\)-norm and become particle columns of \(\mathsf T\); eigenvectors with negative eigenvalues have a negative \(\Sigma_3\)-norm and become hole columns. For an indefinite \(\mathsf H\) this link fails: an eigenvector with a positive eigenvalue can have a negative \(\Sigma_3\)-norm, which is the negative-energy mode of the indefinite case in Paraunitary Diagonalization (Section~\ref{sec-lswt-paraunitary-diagonalization}).

\textbf{Diagonalizability.} The Cholesky factorization \(\mathsf H=\mathsf K\mathsf K^\dagger\) makes \(\Sigma_3\mathsf H=\Sigma_3\mathsf K\mathsf K^\dagger\) similar, through \(\mathsf K^\dagger\), to the Hermitian matrix \(\mathsf K^\dagger\Sigma_3\mathsf K\):

\[
\mathsf K^\dagger\big(\Sigma_3\mathsf H\big)(\mathsf K^\dagger)^{-1}=\mathsf K^\dagger\Sigma_3\mathsf K.
\]

A Hermitian matrix is diagonalizable with real eigenvalues, so \(\Sigma_3\mathsf H\) is diagonalizable with the same eigenvalues. For a positive-semidefinite \(\mathsf H\), Jordan blocks, such as the one of the antiferromagnetic Goldstone mode, can therefore occur only when \(\mathsf H\) is singular. For an indefinite \(\mathsf H\) the Cholesky argument does not apply, and \(\Sigma_3\mathsf H\) can fail to be diagonalizable even when \(\mathsf H\) is invertible. An example with two modes is \(\mathsf A=\operatorname{diag}(1,-0.7)\) and an anomalous coupling \(0.15\) between the two modes, for which \(\det\mathsf H\approx0.522\) while \(\Sigma_3\mathsf H\) has a \(2\times2\) Jordan block at the eigenvalue \(0.85\). Such a case lies outside the stable reference configurations treated by LSWT.

\textbf{Equal numbers of positive and negative eigenvalues.} The matrix \(\mathsf K^\dagger\Sigma_3\mathsf K\) is congruent to \(\Sigma_3\) through the invertible matrix \(\mathsf K\). By Sylvester's law of inertia, congruent Hermitian matrices have the same numbers of positive and negative eigenvalues, so \(\Sigma_3\mathsf H\) has \(N_{\mathrm{sub}}\) positive and \(N_{\mathrm{sub}}\) negative eigenvalues.

\subsection{Existence of the Paraunitary Transformation}

The three properties construct \(\mathsf T\) directly from eigenvectors. Within an eigenspace of \(\Sigma_3\mathsf H\) with eigenvalue \(\lambda\), the form \(\mathbf w^\dagger\Sigma_3\mathbf v\) is definite with the sign of \(\lambda\), because \eqref{eq-lswt-metric-sign-rule} holds for every vector of the eigenspace. Gram--Schmidt orthogonalization with respect to this definite form gives a basis with \(\mathbf v^\dagger\Sigma_3\mathbf v=\operatorname{sign}\lambda\). Eigenvectors of different eigenvalues are already \(\Sigma_3\)-orthogonal by \eqref{eq-lswt-metric-orthogonality}, since all eigenvalues are real. Collecting the \(N_{\mathrm{sub}}\) vectors of positive eigenvalues first and the \(N_{\mathrm{sub}}\) vectors of negative eigenvalues second gives a matrix \(\mathsf T\) with \(\mathsf T^\dagger\Sigma_3\mathsf T=\Sigma_3\). The eigenvalue equation \(\Sigma_3\mathsf H\mathsf T=\mathsf T\,\Sigma_3\mathsf D\), with \(\mathsf D\) the diagonal matrix of the absolute values of the eigenvalues, then gives \(\mathsf T^\dagger\mathsf H\mathsf T=\mathsf T^\dagger\Sigma_3\mathsf T\,\Sigma_3\mathsf D=\mathsf D\), a diagonal matrix with positive entries.

Conversely, if a paraunitary \(\mathsf T\) with \(\mathsf T^\dagger\mathsf H\mathsf T=\mathsf D\) positive and diagonal exists, then \(\mathsf H=(\mathsf T^{-1})^\dagger\mathsf D\,\mathsf T^{-1}\) is congruent to a positive diagonal matrix and is positive definite. A paraunitary transformation to a diagonal form with positive entries therefore exists if and only if \(\mathsf H\) is positive definite.

\subsection{Colpa's Transformation Satisfies Both Conditions}

Colpa's construction, \(\mathsf T=(\mathsf K^\dagger)^{-1}\mathsf V(\Lambda\Sigma_3)^{1/2}\) with \(\mathsf V\Lambda\mathsf V^\dagger=\mathsf K^\dagger\Sigma_3\mathsf K\), is verified in Paraunitary Diagonalization (Section~\ref{sec-lswt-paraunitary-diagonalization}) through \(\mathsf T\Sigma_3\mathsf T^\dagger=\Sigma_3\) and \(\mathsf T^\dagger\mathsf H\mathsf T=\Lambda\Sigma_3\). The second form of the paraunitary condition can also be checked directly. The matrix \(\mathsf K^\dagger\Sigma_3\mathsf K\) is invertible, so \(\Lambda\) is invertible, and \(\Sigma_3^{-1}=\Sigma_3\) gives \(\mathsf K^{-1}\Sigma_3(\mathsf K^\dagger)^{-1}=(\mathsf K^\dagger\Sigma_3\mathsf K)^{-1}=\mathsf V\Lambda^{-1}\mathsf V^\dagger\). Hence

\[
\mathsf T^\dagger\Sigma_3\mathsf T
=(\Lambda\Sigma_3)^{1/2}\mathsf V^\dagger\,\mathsf K^{-1}\Sigma_3(\mathsf K^\dagger)^{-1}\,\mathsf V(\Lambda\Sigma_3)^{1/2}
=(\Lambda\Sigma_3)^{1/2}\Lambda^{-1}(\Lambda\Sigma_3)^{1/2}
=\Sigma_3,
\]

where the last step uses that the diagonal matrices commute and \(\Lambda\Sigma_3\Lambda^{-1}=\Sigma_3\). This is the check left in commented form in the restructured TeX appendix; by the first section it is equivalent to \(\mathsf T\Sigma_3\mathsf T^\dagger=\Sigma_3\).

\subsection{Freedom in the Transformation}

Let \(\mathsf T\) and \(\mathsf T'\) both be paraunitary with \(\mathsf T^\dagger\mathsf H\mathsf T=\mathsf T'^\dagger\mathsf H\mathsf T'=\mathsf D\), with the same ordering of the diagonal entries. The matrix \(\mathsf U=\mathsf T^{-1}\mathsf T'\) is paraunitary by the group property. From \(\Sigma_3\mathsf H\mathsf T=\mathsf T\Sigma_3\mathsf D\) and the same relation for \(\mathsf T'\),

\[
\mathsf U\,\Sigma_3\mathsf D=\mathsf T^{-1}\Sigma_3\mathsf H\mathsf T'=\Sigma_3\mathsf D\,\mathsf U,
\]

so \(\mathsf U\) commutes with the diagonal matrix \(\Sigma_3\mathsf D\) of signed energies and is block diagonal in its eigenspaces. When \(\mathsf H\) is positive definite, each eigenspace contains only particle columns or only hole columns, because their signed energies have opposite signs. On such a block \(\Sigma_3\) is \(+\mathsf I\) or \(-\mathsf I\), and the paraunitary condition restricted to the block is unitarity. The transformation is therefore unique up to a phase for each nondegenerate column and a unitary mixing within each degenerate set of columns.

\subsection{Particle--Hole Structure of the Transformation}

The Bogoliubov transformation has the block form \(\mathsf T_{\mathbf k}=\begin{pmatrix}\mathsf P_{\mathbf k}&\mathsf Q_{-\mathbf k}\\ \mathsf Q_{\mathbf k}^*&\mathsf P_{-\mathbf k}^*\end{pmatrix}\), so that the lower components of \(\hat\Psi_{\mathbf k}\) are the adjoints of the upper components of \(\hat\Psi_{-\mathbf k}\). With \(\Sigma_1=\sigma_1\otimes\mathsf I_{N_{\mathrm{sub}}}\), which exchanges the two block rows when applied on the left and the two block columns when applied on the right, the block form is equivalent to

\begin{equation}\label{eq-lswt-transformation-particle-hole}
\mathsf T_{\mathbf k}=\Sigma_1\mathsf T_{-\mathbf k}^*\Sigma_1.
\end{equation}

A transformation of this form can always be chosen. Let \(\mathbf t\) be a particle column of \(\mathsf T_{-\mathbf k}\) with \(\Sigma_3\mathsf H_{-\mathbf k}\mathbf t=\varepsilon\,\mathbf t\) and \(\mathbf t^\dagger\Sigma_3\mathbf t=1\). The particle--hole relation \(\Sigma_1\mathsf H_{-\mathbf k}^*\Sigma_1=\mathsf H_{\mathbf k}\) and \(\Sigma_3\Sigma_1=-\Sigma_1\Sigma_3\) give \(\Sigma_3\mathsf H_{\mathbf k}(\Sigma_1\mathbf t^*)=-\varepsilon\,\Sigma_1\mathbf t^*\), and

\[
(\Sigma_1\mathbf t^*)^\dagger\Sigma_3(\Sigma_1\mathbf t^*)=\mathbf t^{\mathsf T}\Sigma_1\Sigma_3\Sigma_1\mathbf t^*=-\big(\mathbf t^\dagger\Sigma_3\mathbf t\big)^*=-1.
\]

The vector \(\Sigma_1\mathbf t^*\) is therefore a correctly normalized hole column of \(\mathsf T_{\mathbf k}\) with signed energy \(-\varepsilon\), which is \(-\varepsilon_{n,-\mathbf k}\) for the band of \(\mathbf t\). Building the hole columns of \(\mathsf T_{\mathbf k}\) in this way from the particle columns of \(\mathsf T_{-\mathbf k}\), and doing the same with \(\mathbf k\) and \(-\mathbf k\) exchanged, yields \eqref{eq-lswt-transformation-particle-hole}. The resulting columns are \(\Sigma_3\)-orthogonal to the particle columns of \(\mathsf T_{\mathbf k}\) by \eqref{eq-lswt-metric-orthogonality}, since their signed energies have the opposite sign. A transformation computed independently at \(\mathbf k\), such as Colpa's, has hole columns that agree with these up to the freedom of the previous section. At a momentum with \(-\mathbf k\equiv\mathbf k\) modulo a reciprocal-lattice vector, the relation constrains \(\mathsf T_{\mathbf k}\) itself and fixes part of that freedom.

\subsection{References}

\subsubsection{Internal Documents}

\begin{itemize}
\tightlist
\item
  Paraunitary Diagonalization (Section~\ref{sec-lswt-paraunitary-diagonalization}): states the paraunitary condition, the particle--hole structure of the spectrum, Colpa's construction, and the gauge freedom proved here.
\item
  \lswtnoteref{LSWT boson Hamiltonian}{Momentum-Space BdG Hamiltonian}: defines the BdG blocks and the block relations behind the particle--hole relation.
\item
  \lswtnoteref{LSWT spin Hamiltonian}{Notation and Conventions}: defines the Nambu metric and the paraunitary-matrix notation.
\item
  Worked Example (Section~\ref{sec-lswt-worked-example}): evaluates these statements for a single sublattice.
\end{itemize}

\subsubsection{External Sources}

\begin{itemize}
\tightlist
\item
  J. H. P. Colpa, ``Diagonalization of the Quadratic Boson Hamiltonian,'' \emph{Physica A} \textbf{93}, 327--353 (1978), \href{https://doi.org/10.1016/0378-4371(78)90160-7}{doi:10.1016/0378-4371(78)90160-7}: metric-preserving diagonalization of a positive-definite quadratic boson Hamiltonian through the Cholesky factorization.
\end{itemize}
